\section{Periods in closed form}\label{sec:periods}
\subsection{The normalized theta expression}
We use the odd Jacobi function and the Dedekind function in FAM C.1:
\[
\vartheta_1(v,\tau)=2\sum_{n\ge0}(-1)^n e^{\pi i\tau(n+1/2)^2}\sin((2n+1)\pi v),\qquad
\eta_D(\tau)=e^{\pi i\tau/12}\prod_{n\ge1}(1-e^{2\pi in\tau}).
\]
The cusp branch is reached from $\tau=iT$, $T>1$, on which $\mu=1/2-iy(T)$ and $0<y(T)<T/2$.
Set $a_{\rm cusp}=5+2\sqrt6$. The elliptic-logarithm interpretation is due to \cite[Remark~3.12]{AC26}; the zero formula is due to Eichler and Zagier \cite[unnumbered theorem, p.~399]{EZ82}, also stated in \cite[§2, equation~(6)]{DI10}.

\begin{proposition}[FAM C.1.1; T]\label{fam:C.1.1}
For the period section identified in Statement~\ref{fam:B.1.3}, the particular solution
with \((\beta^{\rm part}+\tau)(p_0)=0\) is
\begin{equation}\label{fam:eq:cd-beta}
 \beta^{\rm part}=\frac{6i}{\pi}\operatorname{Log}
 \frac{\vartheta_1(\mu,\tau)}{2\sqrt3\,\eta_D(\tau)}.
\end{equation}
Here the logarithm is continued on the simply connected uniformizing
half-plane from its cusp expansion. Its special values are
\[
 \mu_0=\frac12-\frac{i\log a_{\rm cusp}}{2\pi},\qquad
 \mu(z_1)=\frac{2-\rho}{3},\qquad \mu(z_2)=\frac{1-i}{2},
\]
\[
 \beta^{\rm part}(z_2)=i\left(\frac32-\frac3\pi\log6\right),
 \qquad \operatorname{Re}\beta^{\rm part}(z_1)=-\frac12.
\]
In particular \(e^{2\pi i\mu_0}=-a_{\rm cusp}\).
\end{proposition}
\noindent\textit{Uses:} FAM B.1.3; the cusp normalization; the Jacobi transformation laws.\par
\begin{proof}
\noindent\textbf{Step 1.} Statement~\ref{fam:B.1.3} identifies the elliptic point represented by \(\mu\) with
one of the sections \((0,\pm\sqrt2t^2(t-1))\). The uniformization map
\((v,\tau)\mapsto(f_{\rm sc}^2\wp_\tau(v),f_{\rm sc}^3\wp'_\tau(v)/2)\) therefore gives
\(\wp_\tau(\mu)=0\). No zero of \(\wp\) is a lattice point. Thus the numerator
of \eqref{fam:eq:cd-beta} is nonzero; the denominator is nonzero as well.
A holomorphic logarithm exists on the simply connected cover.


\noindent\textbf{Step 2.} For completeness the cusp value follows directly from the normally
convergent Fourier expansion, rather than from a quoted Eichler--Zagier
constant:
\begin{equation}\label{fam:eq:cd-wp}
 \frac{\wp_\tau(v)}{\pi^2}
 =\csc^2(\pi v)-\frac13+
 8\sum_{n\geq1}\frac{nq_\tau^n}{1-q_\tau^n}(1-\cos(2\pi nv)),
 \qquad q_\tau=e^{2\pi i\tau}.
\end{equation}
This identity follows by differentiating the logarithm of the Jacobi
product twice and using
\(E_2=1-24\sum nq_\tau^n/(1-q_\tau^n)\). Boundedness of the chosen elliptic logarithm
at the cusp permits passage to the limit. We obtain
\(\sin^2(\pi\mu_0)=3\), equivalently
\(v_{\rm cusp}^2+10v_{\rm cusp}+1=0\), \(v_{\rm cusp}=e^{2\pi i\mu_0}\).
On the chosen imaginary-axis branch \(\operatorname{Im}\mu<0\), so \(v_{\rm cusp}=-a_{\rm cusp}\).
One can see the branch without a numerical convention: on the rectangular
elliptic curve with \(\tau=iT\), \(T>1\), \(\wp\) is real on the segment
\(1/2-iy\), \(0\leq y\leq T/2\). Its endpoint values are the largest and
middle real roots of its cubic, respectively positive and negative.
Its derivative has no zero in the open segment, since the zeros of
\(\wp'\) are the three nonzero half-periods. There is exactly one zero
there; this gives the indicated continuation. Consequently
\(2\sin(\pi\mu_0)=2\sqrt3\).


\noindent\textbf{Step 3.} Put \(R=\vartheta_1/\eta_D\). Poisson summation for the defining Gaussian
series and its index shift give
\[
 \begin{aligned}
 R(v,\tau+1)&=e^{\pi i/6}R(v,\tau),\\
 R(v/\tau,-1/\tau)&=-i e^{\pi iv^2/\tau}R(v,\tau),\\
 R(1-v,\tau)&=R(v,\tau),\qquad R(v+1,\tau)=-R(v,\tau).
 \end{aligned}
\]
Inserting the actual transformed arguments gives
\[
 \begin{aligned}
 R(\mu\circ g_1,\tau\circ g_1)
 &=e^{-\pi i/3}e^{\pi i(1-\mu)^2/\tau}R(\mu,\tau),\\
 R(\mu\circ g_2,\tau\circ g_2)
 &=i e^{\pi i\mu^2/\tau}R(\mu,\tau).
 \end{aligned}
\]
Multiplication of the logarithms by \(6i/\pi\) gives exactly
\(2-6(1-\mu)^2/\tau\) and \(-3-6\mu^2/\tau\), modulo constants in
\(12\mathbb Z\). These constants vanish: each defect is constant on the
connected cover, and evaluation at the corresponding fixed point makes
both its two sides zero. The fixed-point equations themselves give
\(\mu(z_1)=(2-\rho)/3\) and \(\mu(z_2)=(1-i)/2\).
The Jacobi product has cusp expansion
\(R(\mu,\tau)=2q_\tau^{1/12}\sin(\pi\mu)(1+O(q_\tau))\).
Equation \eqref{fam:eq:cd-beta} consequently gives
\(\beta^{\rm part}+\tau\to0\). The difference from the specified particular
solution is an invariant holomorphic function extending across all three
points of \(B\), hence a constant, and the cusp normalization makes it zero.


\noindent\textbf{Step 4.} At \(\tau=i\), the half-period translation formula gives
\(\vartheta_1((1-i)/2,i)=e^{\pi/4}\vartheta_3(0,i)\).
The identities \(2\eta_D^3=\vartheta_2\vartheta_3\vartheta_4\),
\(\vartheta_2(0,i)=\vartheta_4(0,i)=2^{-1/4}\vartheta_3(0,i)\)
give \(R=e^{\pi/4}\sqrt2\), proving the displayed value of \(\beta\).
At \(\rho\), \(q_\tau\) is negative real, \(\mu=1/2-i/(2\sqrt3)\), and the
product for \(R\) has phase \(\pi/12\), with all its remaining factors
positive real. This gives \(\operatorname{Re}\beta=-1/2\).
\end{proof}

\begin{computed}[FAM C.1.2; T]\label{fam:C.1.2}
The following decimal evaluations of the exact expressions in C.1.1 agree
in two independent evaluations through the displayed twenty decimal places:
\[
\begin{array}{c|l}
\mu_0&0.5-0.36485183191106784792\,i\\
\mu(z_1)&0.5-0.28867513459481288225\,i\\
\mu(z_2)&0.5-0.5\,i\\
\beta^{\rm part}(z_1)&-0.5-1.27101206057537346287\,i\\
\beta^{\rm part}(z_2)&-0.21100425815613411418\,i
\end{array}
\]
These are evaluations of explicit convergent expressions, with working
precision ninety decimal digits. The formula \eqref{fam:eq:cd-beta}, together
with \(j(\tau)=1728t\) and \(\wp_\tau(\mu)=0\), determines the periods
throughout a fundamental domain.
\end{computed}
\noindent\textit{Uses:} FAM C.1.1; the theta product and its displayed tail bound.\par
\begin{proof}
\noindent\textbf{Step 1.} The source run records that the executable \texttt{cd\_numerics.py} evaluates \(\mu_0\) first by its
algebraic logarithm and then by solving
\(\pi^2\csc^2(\pi v)-\pi^2/3=0\). For \(\beta\), its first calculation uses
the Jacobi theta series and the eta product; its second uses
\[
 R(v,\tau)=2e^{\pi i\tau/6}\sin(\pi v)
 \prod_{n\geq1}(1-2q_\tau^n\cos(2\pi v)+q_\tau^{2n}).
\]
The product is truncated after 120 factors. At the two specified points
\(|q_\tau|\leq e^{-\pi\sqrt3}\) and \(|\operatorname{Im}v|\leq1/2\), so the
omitted logarithm is bounded by
\(\sum_{n>120}(2|q_\tau|^n\cosh\pi+|q_\tau|^{2n})/(1-2|q_\tau|^n\cosh\pi-|q_\tau|^{2n})\),
which is less than \(10^{-280}\). The measured discrepancies of the
ninety-digit calculations are below \(10^{-88}\). We do not promote floating-point rounding to an interval certificate.


\noindent\textbf{Step 2.} 
A further sixty-digit interval computation, \texttt{cd\_special\_intervals.py},
evaluates the real positive product at $\rho$, including its forty-factor
tail, and the exact logarithmic expression at $i$. It encloses both
imaginary parts within intervals narrower than $3\cdot10^{-59}$ and
therefore certifies the displayed twenty decimal places. It also proves
$-0.2<\operatorname{Im}(\kappa_1^{\rm part}-\kappa_2^{\rm part})<-0.1$.
The exact formulas, product-tail bounds, independent algorithms, and
full outputs are reported by the source run as accompanying the decimals. The special values of \(\mu\) are also checked directly in
the fixed-point equations, an exact algebraic calculation.
\end{proof}

The computation in FAM C.1.2 is the recorded source computation. Neither \path{cd_numerics.py} nor \path{cd_special_intervals.py}, nor their full outputs, is among the supplied attachments; no current execution is claimed. The latter interval is the input used by FAM A.4.1 for its generic Albanese cross-check. All program and output-file references in the period and maximum proofs below likewise describe the source runs; they do not assert a current execution or the presence of those artifacts in the supplied attachments.

\subsection{The elliptic norm and its special values}
The elliptic lattice in this subsection is denoted $\Lambda_{\rm ell}$ to distinguish it from the marked rank-four lattice. Put
\[
\mathcal N(v,\tau)=e^{-\pi(\operatorname{Im}v)^2/\operatorname{Im}\tau}
\left|\frac{\vartheta_1(v,\tau)}{\eta_D(\tau)}\right|,
\qquad U_\tau(v)=\log\mathcal N(v,\tau),
\qquad \mathcal L(t)=U_\tau(\mu).
\]
For $\Lambda_{\rm ell}=\omega(\mathbb Z+\tau\mathbb Z)$ use $U_{\Lambda_{\rm ell}}(z)=U_\tau(z/\omega)$.
The theta formula gives the source identity (FAM C.2.1, equation~\texttt{cd-norm})
\begin{equation}\label{fam:eq:cd-norm}
D_{\rm part}=\frac6\pi\log\frac{\mathcal N(\mu,\tau)}{2\sqrt3}.
\end{equation}

\begin{proposition}[ARI C.1.1; T; distribution under isogenies]\label{ari:C.1.1}
Let $\Lambda_{\mathrm{ell}},\Lambda_{\mathrm{ell}}'\subset\mathbb C$ be lattices and let
$\alpha\ne0$ satisfy $\alpha\Lambda_{\mathrm{ell}}\subset\Lambda_{\mathrm{ell}}'$. Write
$\phi:\mathbb C/\Lambda_{\mathrm{ell}}\to\mathbb C/\Lambda_{\mathrm{ell}}'$ for $z\mapsto\alpha z$,
$K=\alpha^{-1}\Lambda_{\mathrm{ell}}'/\Lambda_{\mathrm{ell}}$, and
$d=[\Lambda_{\mathrm{ell}}':\alpha\Lambda_{\mathrm{ell}}]$. Then
\begin{equation}\label{ari:eq:cs-isogeny}
 U_{\Lambda_{\mathrm{ell}}'}(\alpha z)=\sum_{v\in K}U_{\Lambda_{\mathrm{ell}}}(z+v).
\end{equation}
The equality holds away from the zeros and as an equality of locally
integrable functions at them. If $\Delta_{\Lambda_{\mathrm{ell}}}=g_2(\Lambda_{\mathrm{ell}})^3-
27g_3(\Lambda_{\mathrm{ell}})^2$, then
\[
 \prod_{0\ne v\in K}\mathcal N_{\Lambda_{\mathrm{ell}}}(v)
 =|\alpha|\left|\frac{\Delta_{\Lambda_{\mathrm{ell}}'}}{\Delta_{\Lambda_{\mathrm{ell}}}}\right|^{1/12}.
\]
In particular, for every nonzero endomorphism of one elliptic curve,
$\prod_{v\in\ker\alpha,\ v\ne0}\mathcal N(v)=|\alpha|$.
\end{proposition}
\noindent\textit{Uses:} the theta product and Haar probability measure.\par
\begin{proof}
\noindent\textbf{Step 1.} The Jacobi product, with $q_\tau=e^{2\pi i\tau}$ and
$X=e^{2\pi iz}$, gives
\[
 \vartheta _1(z,\tau)/\eta_D(\tau)
 =i q_\tau^{1/12}X^{-1/2}(1-X)
     \prod_{m\ge1}(1-q_\tau^mX)(1-q_\tau^mX^{-1}),
\]
where a harmless fixed unit changes with the theta convention. For
$0<y<T$, integration of the logarithm over $z=x+iy$, $0\le x\le1$,
gives $-\pi T/6+\pi y-\pi y^2/T$. Each product factor has mean logarithm
zero by the absolutely convergent logarithmic series; the endpoint
singularities are integrable. Integrating in $y$ proves
$\int U_{\Lambda_{\mathrm{ell}}}\,d\mu_{\Lambda_{\mathrm{ell}}}=0$ for Haar probability measure.
Moreover
\[
 \Delta_{\rm Eucl}U_{\Lambda_{\mathrm{ell}}}
   =2\pi\delta_0-\frac{2\pi}{\operatorname{area}\Lambda_{\mathrm{ell}}},\qquad
 U_{\Lambda_{\mathrm{ell}}}(z)=\log|z|+\frac1{12}\log|\Delta_{\Lambda_{\mathrm{ell}}}|+o(1).
\]
The second formula follows from
$\vartheta _1'(0,\tau)=2\pi\eta_D(\tau)^3$ and
$\Delta_{\mathbb Z+\tau\mathbb Z}=(2\pi)^{12}\eta_D(\tau)^{24}$,
then scaling by $\omega$. The two sides of
\eqref{ari:eq:cs-isogeny} have the same logarithmic masses at $K$ and the
same constant Laplacian, since
$d=|\alpha|^2\operatorname{area}\Lambda_{\mathrm{ell}}/
\operatorname{area}\Lambda_{\mathrm{ell}}'$. Their difference is a smooth harmonic
function on the compact torus, hence constant. Both means vanish:
translations preserve Haar measure and an isogeny sends Haar
probability measure to Haar probability measure. The constant is zero.
Taking $z\to0$ and cancelling the logarithmic singularity gives the
product. The kernel, Laplacian and mean-zero argument fix the normalization.
\end{proof}

\begin{proposition}[ARI C.1.2; T; the phases and the two periods]\label{ari:C.1.2}
With the branch continued from the specified cusp, the exact values are
\begin{align*}
 \frac{\vartheta _1((2-\rho)/3,\rho)}{\eta_D(\rho)}
   &=e^{\pi i/12}\,3^{1/4}e^{\pi/(6\sqrt3)},\\
 \frac{\vartheta _1((1-i)/2,i)}{\eta_D(i)}
   &=\sqrt2\,e^{\pi/4},\\
 \beta^{\rm part}(z_1)
   &=-\frac12+i\left(\frac1{\sqrt3}-\frac3{2\pi}\log48\right),\\
 \beta^{\rm part}(z_2)
   &=i\left(\frac32-\frac3\pi\log6\right).
\end{align*}
Consequently $\mathcal N(z_1)=3^{1/4}$, $\mathcal N(z_2)=\sqrt2$,
$D_{\rm part}(z_1)=-3\log48/(2\pi)$, and
$D_{\rm part}(z_2)=-3\log6/\pi$.
\end{proposition}
\noindent\textit{Uses:} ARI C.1.1; FAM C.1.1; the fixed cusp branch.\par
\begin{proof}
\noindent\textbf{Step 1.} By ARI C.1.1, multiplication by $1+\rho$ has matrix
$\left(\begin{smallmatrix}1&-1\\1&2\end{smallmatrix}\right)$,
its kernel is $\{0,\mu,-\mu\}$, and $(1+\rho)\mu=1$ for
$\mu=(2-\rho)/3$. Statement ARI C.1.1 gives $\mathcal N(\mu)^2=\sqrt3$.
For the requested parallel calculation at $i$, multiplication by $1+i$
has kernel $\{0,(1-i)/2\}$ and absolute value $\sqrt2$; the same
identity gives $\mathcal N((1-i)/2,i)=\sqrt2$ without a separate theta-constant
evaluation.


\noindent\textbf{Step 2.} At $\rho$, $\mu=1/2-i/(2\sqrt3)$ and $q_\tau=-e^{-\pi\sqrt3}$. In
\[
 R(z,\tau)=2e^{\pi i\tau/6}\sin(\pi z)
 \prod_{m\ge1}\bigl(1-2q_\tau^m\cos(2\pi z)+q_\tau^{2m}\bigr),
\]
$\sin(\pi\mu)$ is positive real. The factors are
$(1+q_\tau^m e^{\pi/\sqrt3})(1+q_\tau^m e^{-\pi/\sqrt3})$, both positive:
when $m$ is odd the larger subtracted number is at most
$e^{-2\pi/\sqrt3}<1$. Thus the phase is exactly $\pi/12$ modulo
$2\pi$. There is no winding ambiguity in the continued logarithm.
Indeed continue from the cusp along $\tau=1/2+iT$,
$T\ge\sqrt3/2$. Real conjugation of the rhombic lattice keeps the
chosen zero on $\mu=1/2-iy$, $0<y<T$. It cannot leave that component:
a collision of the two zeros would require $g_3=0$, which does not
occur on this ray. Indeed, writing $q_\tau=-r$, $0<r<1/200$,
the tail after the first term of the $E_6$ series has absolute value
at most $504\cdot32r^2/((1-r^2)(1-8r))<1/2$, while its
first term is positive. Thus $E_6>1/2$. Hitting $y=0$ would again give a half-period zero, and
hitting $y=T$ would give a pole. Each displayed product factor remains
positive because $e^{-2\pi T}e^{2\pi y}<1$. The cusp phase tends to
$\pi/12$, so the continued argument has that value. On the imaginary
ray $\tau=iT$, $T\ge1$, the same branch is
$\mu=1/2-iy$ with $0<y\le T/2$; every product factor is positive,
so its argument at $i$ is zero. These checks supply the branch step
used in FAM C.1.1. The Gaussian factors are respectively
$e^{\pi/(6\sqrt3)}$ and $e^{\pi/4}$. Substitution in the definition of
$\beta^{\rm part}$ proves all four identities.
\end{proof}

\subsection{The Albanese nomes}

\begin{proposition}[ARI C.2.1; T; exact nomes and continuation radius]\label{ari:C.2.1}
For $w=e^{-\pi ic}$, the fixed branches give
\begin{align*}
 \kappa_1&=c+\tfrac12-\frac{3i}{2\pi}\log48,&
 q_1&=-w^2/48^3,\\
 \kappa_2&=c+\tfrac32-\frac{3i}{\pi}\log6,&
 q_2&=iw/216,\qquad q_1=(27/64)q_2^2.
\end{align*}
The family is parameterized by $0<|q_2|<1$ and exists precisely in
that domain. Writing $J(q_\tau)=q_\tau^{-1}+744+\cdots$, one has
\[
 j(\operatorname{Alb}S_1)=J(-w^2/48^3).
\]
This function is holomorphic on $0<|w|<48^{3/2}$, has a pole of order
two at zero, and has a natural boundary at $|w|=48^{3/2}$.
In particular it continues holomorphically through every point of the
family boundary $|w|=216$.
\end{proposition}
\noindent\textit{Uses:} FAM A.4.1 and A.5.1; ARI C.1.2; the maximum theorem stated in the Result.\par
\begin{proof}
\noindent\textbf{Step 1.} The precise Albanese lattices in FAM A.4.1 and A.5.1 are verified by projecting all
translations on the product covers. At the order-three fibre the
projection is $(1/3)(\mathbb Z+\kappa_1\mathbb Z)$; multiplication by
three gives the first stated lattice. At the order-four fibre it is
$\tfrac12\mathbb Z+\tfrac14\kappa_2\mathbb Z$; multiplication by two
gives $\mathbb Z+(\kappa_2/2)\mathbb Z$. Thus the positive-imaginary
period representatives are $-\kappa_1$ and $-\kappa_2/2$.
Adding $2\mu_1=1-i/\sqrt3$ and
$3\mu_2=3/2-3i/2$ to C.1.2 proves the two $\kappa$ identities,
including their real parts. Exponentiation gives the signs in $q_1,q_2$;
the ratio is $216^2/48^3=27/64$.


\noindent\textbf{Step 2.} The factor $1/2$ in $q_2=iw/216$ is exactly the Albanese projection just computed. The natural-boundary argument uses modular invariance at rational boundary points. The existence criterion and exact $2\mathbb Z$
identification are FAM C.2.3 and FAM A.4.1 and A.5.1; the latter is also Theorem~1.4
(label \texttt{p2:thm:exact-period}) of
\cite[Theorem~1.4 (\texttt{p2:thm:exact-period})]{DGKXv1}. Its hypotheses specify precisely the
retained period family, affine finite twists and cusp collar. They
coincide with the family in this prompt; the normalization differs
only by the fixed choice of the particular solution, which does not
alter parameter differences. Version~2 \cite[Introduction]{DGKXv2} omits that theorem; we infer
no mathematical conclusion from this omission.


\noindent\textbf{Step 3.} The Fourier expansion for $J$ converges for $0<|q_\tau|<1$ and has a simple
pole at $q_\tau=0$. The inverse image under $q_\tau=-w^2/48^3$ is the claimed
punctured disc and the pole has order two. The inequality
$216<48^{3/2}$ is equivalent to $27<64$. At every nonzero point on the
larger boundary the map to $q_\tau$ is locally biholomorphic. A continuation
there would continue $J$ through the unit circle, contradicting the
same rational-cusp argument. The phrase ``up to the
radius'' therefore means the open punctured disc, not holomorphy on
its bounding circle or at its central pole.
\end{proof}

\subsection{Division values}
For division-polynomial calculations we use the rational model of ARI C.4:
\[
E^{\rm div}_t:\quad y_{\rm div}^2=x_{\rm div}^3+A(t)x_{\rm div}+B(t),\quad
A=-\tfrac34t^3(t-1),\quad B=\tfrac14t^4(t-1)^2,
\]
\[
P^{\rm div}=(0,t^2(t-1)/2),\qquad \Delta_{\rm div}(t)=27t^8(t-1)^3.
\]
Its coordinates are distinguished from the cubic coordinates in FAM B.1.1. The notation sheet records the homothety and the sign of its distinguished section; the invariant norm is the same. We use $\lambda_\infty(Q)=-\log\mathcal N(Q)$ as the normalized local-height notation of ARI C.4.1.

\begin{proposition}[ARI C.4.2; T; the full division identity]\label{ari:C.4.2}
Let $\psi_m$ be the standard division polynomial of a smooth
Weierstrass model, normalized by $\psi_1=1$, $\psi_2=2y$.
For $m\ge2$ and $Q,mQ\ne O$,
\begin{equation}\label{ari:eq:cs-division}
 \mathcal N(Q)^{m^2}
       =\mathcal N(mQ)\frac{|\Delta_E|^{(m^2-1)/12}}{|\psi_m(Q)|}.
\end{equation}
For the displayed family this specializes to
\[
 \mathcal N(P^{\mathrm{div}})^4=3^{3/4}|t-1|^{-1/4}\mathcal N(2P^{\mathrm{div}}),\qquad
 \mathcal N(P^{\mathrm{div}})^9=16|t|^{-2/3}\mathcal N(3P^{\mathrm{div}}).
\]
If $P^{\mathrm{div}}_t$ has exact order $n$ and $m\ge2$ satisfies
$m\equiv1$ or $-1\pmod n$, then
\begin{equation}\label{ari:eq:cs-torsionvalue}
 \mathcal N(P^{\mathrm{div}}_t)^{12}
       =\frac{|\Delta_{\mathrm{div}}(t)|}{|\psi_m(P^{\mathrm{div}}_t)|^{12/(m^2-1)}}.
\end{equation}
The positive value on the right is independent of $m$.
\end{proposition}
\noindent\textit{Uses:} ARI C.1.1; the displayed rational Weierstrass model and standard division-polynomial normalization.\par
\begin{proof}
\noindent\textbf{Step 1.} Use the complex uniformization $x=\wp_{\Lambda_{\mathrm{ell}}}(z)$,
$y=\wp'_{\Lambda_{\mathrm{ell}}}(z)/2$. The rational function $\psi_m$ has divisor
$\sum_{0\ne Q\in E[m]}(Q)-(m^2-1)(O)$ and leading Laurent term
$(-1)^{m-1}m z^{1-m^2}$. One way to verify both normalizations without
a model convention is to form
$(-1)^{m-1}\sigma_{\Lambda_{\mathrm{ell}}}(mz)/\sigma_{\Lambda_{\mathrm{ell}}}(z)^{m^2}$.
Its lattice factors cancel, its zeros are exactly the nonzero
$m$-division points, and its expansion at zero is the displayed
one. The addition formula for $\sigma$ gives the standard
division-polynomial recurrence with initial values $1,2y$, so it is
the stated rational function.


\noindent\textbf{Step 2.} The function
\[
 m^2U_{\Lambda_{\mathrm{ell}}}(z)-U_{\Lambda_{\mathrm{ell}}}(mz)
                    +\log|\psi_m(Q(z))|
\]
is harmonic after its logarithmic singularities cancel. Its
Laplacian has zero constant term as well, so it is constant.
The expansion in ARI C.1.1 and the leading term of $\psi_m$ make that
constant $(m^2-1)\log|\Delta_E|/12$. This proves
\eqref{ari:eq:cs-division}. Direct substitution of
\[
 \psi_2(P^{\mathrm{div}})=t^2(t-1),\qquad
 \psi_3(P^{\mathrm{div}})=-9t^6(t-1)^2/16
\]
gives the two special identities. For $m\equiv\pm1\pmod n$,
$mP^{\mathrm{div}}=\pm P^{\mathrm{div}}$ and $\mathcal N$ is even. Cancelling $\mathcal N(P^{\mathrm{div}})$ in
\eqref{ari:eq:cs-division} gives \eqref{ari:eq:cs-torsionvalue}.
Every right side is the same positive number $\mathcal N(P^{\mathrm{div}})^{12}$, which
proves independence of $m$, not merely numerical agreement.
\end{proof}

\begin{proposition}[ARI C.4.3; T; algebraicity and explicit fields for every order]\label{ari:C.4.3}
Let $t_0\in\mathbb C\setminus\{0,1\}$ and suppose $P^{\mathrm{div}}_{t_0}$ has
exact order $n\ge2$. Then $t_0$ and $\mathcal N(P^{\mathrm{div}}_{t_0})$ are algebraic.
Put $K_0=\mathbb Q(t_0,\bar t_0)$,
\[
 W_{\rm rad}=\Delta_{\mathrm{div}}(t_0)\Delta_{\mathrm{div}}(\bar t_0)>0,\qquad
 V_{\rm rad}=\psi_{n+1}(P^{\mathrm{div}}_{t_0})\psi_{n+1}(P^{\mathrm{div}}_{\bar t_0})>0.
\]
With positive real radicals,
\[
 \mathcal N(P^{\mathrm{div}}_{t_0})^{12n}=W_{\rm rad}^{n/2}V_{\rm rad}^{-6/(n+2)}
 \ \in\ K_0\bigl(W_{\rm rad}^{1/2},V_{\rm rad}^{1/(n+2)}\bigr).
\]
This supplies an explicit field, without a claim that it is minimal.
More generally, if $m\equiv\pm1\pmod n$, then
\[
 \lambda_\infty(P^{\mathrm{div}}_{t_0})
 =\frac{\log|\psi_m(P^{\mathrm{div}}_{t_0})|}{m^2-1}
       -\frac1{12}\log|\Delta_{\mathrm{div}}(t_0)|.
\]
These formulas determine the height and $\mathcal N^{12}$ on the entire
torsion locus.
\end{proposition}
\noindent\textit{Uses:} ARI C.4.2; the nonreal cusp value of the chosen elliptic logarithm.\par
\begin{proof}
\noindent\textbf{Step 1.} The division polynomial $\psi_n(P^{\mathrm{div}}_t)$ is a nonzero rational
polynomial in $t$. To check nonvanishing for every $n$, a hypothetical
torsion identity would express the elliptic logarithm as
$a\tau+b$ with constant rational $a,b$. Its bounded cusp limit
forces $a=0$, but that limit has nonzero imaginary part, a
contradiction. Thus $t_0$ is a zero of a nonzero polynomial over
$\mathbb Q$. Formula \eqref{ari:eq:cs-torsionvalue} shows algebraicity
of $\mathcal N$ using only positive rational powers of nonzero algebraic
absolute values. Taking $m=n+1$, raising to the $n$th power, and
using $|\Delta|=W_{\rm rad}^{1/2}$, $|\psi_{n+1}|=V_{\rm rad}^{1/2}$ gives the exact
field statement. The local-height formula is its logarithm.


\noindent\textbf{Step 2.} Here is a polynomial relation involving only $t$ and $r=\mathcal N^{12}$,
if one prefers that form to positive radicals. Remove the factors
at $0,1$ from $\psi_n(P^{\mathrm{div}}_t)$ and divide successively by the
corresponding polynomials for proper divisors of $n$; let $F_n$ be
the resulting primitive polynomial. At a point of exact order
$d\mid n$, the formal group gives
$z([n/d]Q)=(n/d)z(Q)+O(z(Q)^2)$, so the vanishing multiplicities for
$\psi_d$ and $\psi_n$ agree. Hence this cancellation leaves exactly
the exact-order-$n$ divisor, with its multiplicities. Set
$m=n+1$ and $e=m^2-1$. Clearing rational denominators in
\[
 \operatorname{Res}_S\left(
 F_n(S),\
 r^{2e}\{\psi_m(P^{\mathrm{div}}_T)\psi_m(P^{\mathrm{div}}_S)\}^{12}
              -\{\Delta_{\mathrm{div}}(T)\Delta_{\mathrm{div}}(S)\}^{e}\right)
\]
gives a polynomial in $\mathbb Q[T,r]$ that vanishes at
$(t_0,\mathcal N(P^{\mathrm{div}}_{t_0})^{12})$, since $S=\bar t_0$ is a common root.
It is nonzero: at an exact-order-$n$ root no $\psi_m$ vanishes,
so the leading coefficient in $r$ of the resultant is nonzero.
This is an explicit recurrence-and-resultant prescription for
every $n$, including nonreal torsion parameters.
\end{proof}

\begin{proposition}[ARI C.4.4; T; the five values and all orders through eight]\label{ari:C.4.4}
The exact special values are
\[
 \begin{array}{c|c|c}
 t&\operatorname{ord}(P^{\mathrm{div}})&\mathcal N(P^{\mathrm{div}})^{12}\\ \hline
 0&3&27\\
 1&2&64\\
 32/5&4&10\\
 64(47-27\sqrt3)/11&6&17+10\sqrt3\\
 64(47+27\sqrt3)/11&6&10\sqrt3-17 .
 \end{array}
\]
At $t=0,1$ the orders refer to the smooth models after the specified
base changes. On the smooth base, $F_2=F_3=1$ and
\[
 \begin{aligned}
 F_4&=5t-32,\\
 F_5&=89t^2+4736t-4096,\\
 F_6&=11t^2-6016t+8192,\\
 F_7&=881t^4+4745344t^3-12079104t^2\\
    &\hspace{13mm}+24641536t-16777216,\\
 F_8&=49t^4-2614336t^3-7864320t^2\\
    &\hspace{13mm}+19398656t-8388608.
 \end{aligned}
\]
All these polynomials are squarefree. At each root of exact order
$n=4,5,6,7,8$, the positive number $r=\mathcal N^{12}$ is respectively
determined by
\[
 \begin{array}{c|l}
 n&\text{equation for }r\\ \hline
 4&r=64/|t|\\
 5&r^5=432^5/\bigl((1-t)(5t-32)^4\bigr)\\
 6&r^2=1728^2/F_5(t)\\
 7&r^{35}=2^{192}3^{81}/
       \bigl(|t|^8|t-1|^3|F_6(t)|^{12}\bigr)\\
 8&r^4=1728^4/|F_7(t)|.
 \end{array}
\]
The order-five roots are real and below one; the order-six roots
are real and above one. The two complex order-seven roots are
included in the formula.
\end{proposition}
\noindent\textit{Uses:} ARI C.1.2, C.4.2 and C.4.3; the division-polynomial recurrence.\par
\begin{proof}
\noindent\textbf{Step 1.} The division-polynomial recurrence, evaluated at $x=0$,
$y=t^2(t-1)/2$, gives
\[
 \begin{aligned}
 \psi_4(P^{\mathrm{div}})&=-t^{10}(t-1)^4F_4/32,\\
 \psi_5(P^{\mathrm{div}})&=t^{16}(t-1)^6F_5/4096,\\
 \psi_6(P^{\mathrm{div}})&=9t^{24}(t-1)^9F_6/65536,\\
 \psi_7(P^{\mathrm{div}})&=-t^{32}(t-1)^{12}F_7/16777216,\\
 \psi_8(P^{\mathrm{div}})&=-t^{42}(t-1)^{16}F_4F_8/268435456.
 \end{aligned}
\]
The recursion and all coefficients are checked over $\mathbb Q[t]$
by \path{Cspecial_division.py}. The same program computes
$\gcd(F_n,F_n')=1$ and excludes common factors for proper orders.
Thus no point or multiplicity is omitted. Applying
\eqref{ari:eq:cs-torsionvalue} with $m=n-1$ gives every displayed
formula. The signs for $F_5,F_6$ follow directly from their quadratic
formula; in particular $F_5$ is positive at the two roots of $F_6$.
At order six, substituting
$t=64(47\mp27\sqrt3)/11$ gives the exact identities
\[
 F_5(t)(17\pm10\sqrt3)^2=1728^2.
\]
The positive roots are $17+10\sqrt3$ and $10\sqrt3-17$; their
product is $11$. The substitution is checked in
\path{Cspecial_relations.py}. The order-four value follows from
$64/(32/5)=10$, and the two central values are C.1.2.


\noindent\textbf{Step 2.} For clarity, the recurrence used by the exact program is
\[
 \begin{aligned}
 \psi_{2m+1}&=\psi_{m+2}\psi_m^3-\psi_{m-1}\psi_{m+1}^3,\\
 \psi_{2m}&=\frac{\psi_m}{\psi_2}
   (\psi_{m+2}\psi_{m-1}^2-\psi_{m-2}\psi_{m+1}^2).
 \end{aligned}
\]
Its initial values at $P^{\mathrm{div}}$ are $\psi_0=0,\psi_1=1,\psi_2=2y$,
$\psi_3=-A^2$, $\psi_4=4y(-8B^2-A^3)$. No approximate algebra
enters this calculation. The same code records $\psi_9$, and the
separate file \texttt{Cspecial\_psi16.out} records the factored
$\psi_{16}$ for the duplication ladder. A rational polynomial
relation for the nonreal order-seven values is explicitly computed
by \path{Cspecial_relations.py}, by the resultant prescription
of C.4.3.
\end{proof}

\subsection{Integral identities}

\begin{proposition}[ARI C.8.6; T; the Mahler-measure identities]\label{ari:C.8.6}
Let $\Lambda_{\mathrm{ell}}\subset\C$ be a lattice of area $A$, let $v$ be a zero
of $\wp_{\Lambda_{\mathrm{ell}}}$, and let $d\mu=A^{-1}dA_z$ be Haar probability
measure on $\C/\Lambda_{\mathrm{ell}}$. Put $U(z)=\log \mathcal N(z)$ and
$\mathcal L=U(v)$. Then
\begin{align}
 6\mathcal L&=\int_{\C/\Lambda_{\mathrm{ell}}}\log\left|\frac{\wp'(z)^2}{\wp(z)^3}\right|d\mu(z),
 \label{ari:eq:ca-mahler}\\
 \mathcal N(v)^2&=|\Delta(\Lambda_{\mathrm{ell}})|^{1/6}
      \exp\left(-\int_{\C/\Lambda_{\mathrm{ell}}}\log|\wp(z)|d\mu(z)\right).
 \label{ari:eq:ca-wpmean}
\end{align}
If $C(v_{\wp})=1-(g_2/4)v_{\wp}^2-(g_3/4)v_{\wp}^3$, the pushforward of $d\mu$
under $v_{\wp}=1/\wp(z)$ is
\begin{equation}\label{ari:eq:ca-nu}
 d\nu(v_{\wp})=\frac1{2A}\frac{dA_{v_{\wp}}}{|v_{\wp}|\,|C(v_{\wp})|}.
\end{equation}
In particular its unnormalized integral is exactly $2A$, and
\[
 \mathcal N(v)^6=4\exp\left(\int_{\C}\log|C(v_{\wp})|d\nu(v_{\wp})\right).
\]
Writing $e_1,e_2,e_3$ for the three branch values and
$p(x)=\int\log|x-\wp(z)|d\mu(z)$, the maximum theorem is equivalently
\[
 \int\log|C|d\nu\leq\log2,\qquad
 p(0)\geq\frac{p(e_1)+p(e_2)+p(e_3)}3-\frac{\log2}3.
\]
Equality in either inequality occurs exactly when $g_3=0$.
\end{proposition}
\noindent\textit{Uses:} ARI C.1.1; the maximum theorem stated in the Result.\par
\begin{proof}
\noindent\textbf{Step 1.} The mean of $U$ is zero by the theta product calculation in C.1,
and near zero
$U(z)=\log|z|+\frac1{12}\log|\Delta|+o(1)$.
The divisor of $\wp$ is $(v)+(-v)-2(O)$, with multiplicity
two when $v=-v$. Subtract its logarithmic divisor from the
corresponding combination of $U$'s. The result is harmonic
on the compact torus and hence constant. Its expansion at
zero gives
\[
 \log|\wp(z)|=U(z-v)+U(z+v)-2U(z)
                  +\frac16\log|\Delta|-2\mathcal L.
\]
Integration proves \eqref{ari:eq:ca-wpmean}.
Similarly the divisor of $\wp'$ is the sum of the three
nonzero half-periods minus $3(O)$. The multiplication-by-two
identity of C.1 gives the product of $\mathcal N$ at those half-periods
as $2$. Since $\wp'(z)\sim-2z^{-3}$, comparison at zero yields
\[
 \log|\wp'(z)|=\sum_{w\in E[2]\setminus\{O\}}U(z-w)-3U(z)
                         +\frac14\log|\Delta|.
\]
Its mean is $\frac14\log|\Delta|$. Twice this identity minus
three times the preceding mean proves \eqref{ari:eq:ca-mahler}.
All logarithmic singularities are locally integrable, so these
integrations also cover a zero at a half-period.


\noindent\textbf{Step 2.} The map $z\mapsto x=\wp(z)$ has degree two. Its two inverse
sheets each have Jacobian $|\wp'(z)|^{-2}$; branch values form
a measure-zero set. Thus its pushforward density is
$2/(A|4x^3-g_2x-g_3|)$. Under $x=1/v_{\wp}$ the Jacobian is
$|v_{\wp}|^{-4}$ and the denominator is $4|v_{\wp}|^{-3}|C(v_{\wp})|$.
This proves \eqref{ari:eq:ca-nu} and its normalization. The identity
$\wp'^2/\wp^3=4C(v_{\wp})$ gives the next formula.


\noindent\textbf{Step 3.} Replacing $v$ above by a half-period gives
$p(e_j)=\frac16\log|\Delta|-2U(w_j)$; summing and using
$\sum_jU(w_j)=\log2$ proves the centroid formulation.
The equality characterization is that of C.8.5,
since $j=1728$ is equivalent to $g_3=0$ for a nonsingular lattice.
\end{proof}

\begin{proposition}[ARI C.8.7; T; the two complex beta integrals]\label{ari:C.8.7}
For the square and equianharmonic lattices respectively, the
integrals in C.8.6 are
\[
 \int\log|C|\,d\nu=\log2,\qquad
 \int\log|C|\,d\nu=\frac32\log3-2\log2.
\]
They recover $\mathcal N^2=2$ and $\mathcal N^4=3$ with the same normalization.
\end{proposition}
\noindent\textit{Uses:} ARI C.8.6; the gamma duplication and triplication formulas.\par
\begin{proof}
\noindent\textbf{Step 1.} By homogeneity it suffices to take $C(v_{\wp})=1-v_{\wp}^m$, with $m=2$
or $m=3$. For $-1<\Re s<1-1/m$ put
\[
 I_m(s)=\int_{\C} |v_{\wp}|^{-1}|1-v_{\wp}^m|^{s-1}dA_{v_{\wp}}.
\]
The substitution $w=v_{\wp}^m$, with all $m$ sheets counted, gives
\[
 I_m(s)=\frac1m B_{\C}\left(\frac1{2m},\frac{s+1}2\right),
 \quad
 B_{\C}(a,b)=\pi\frac{\Gamma(a)\Gamma(b)\Gamma(1-a-b)}
                   {\Gamma(1-a)\Gamma(1-b)\Gamma(a+b)}.
\]
We include a proof of the integral evaluation. For
$0<a,b$ and $a+b<1$, insert
\[
 |w|^{2a-2}=\frac1{\Gamma(1-a)}
       \int_0^\infty t^{-a}e^{-t|w|^2}\,dt
\]
and the corresponding formula for $|1-w|^{2b-2}$.
Tonelli's theorem applies to the positive integrand. The
Gaussian integral over $w$ is
$\pi(t+s)^{-1}\exp(-ts/(t+s))$.
Putting $t=rv$, $s=r(1-v)$ reduces the remaining integrals
to $\Gamma(1-a-b)$ and the beta integral $B(a,b)$, proving
the displayed formula. The same integrable bounds justify
differentiation in a neighbourhood of $s=0$.


\noindent\textbf{Step 2.} Consequently, with $a=1/(2m)$ and $\psi=\Gamma'/\Gamma$,
\[
 \frac{I_m'(0)}{I_m(0)}
 =\frac12\left(2\psi(1/2)-\psi(1/2-a)-\psi(1/2+a)\right).
\]
The gamma duplication and triplication formulas, differentiated
logarithmically, evaluate this as $\log2$ for $m=2$ and
$\frac32\log3-2\log2$ for $m=3$. These can also be checked
directly from
$\psi(1/2)=-\gamma-2\log2$,
$\psi(1/4)+\psi(3/4)=-2\gamma-6\log2$ and
$\psi(1/3)+\psi(2/3)=-2\gamma-3\log3$.
Equation \eqref{ari:eq:ca-mahler} completes the normalization test.
\end{proof}

\subsection{Announcements on the norm and the Betti measure}

\begin{announcement}[ARI C.8.8; A; the signed Gaussian identity]\label{ari:C.8.8}
For $\Lambda_{\mathrm{ell}}=\Z+\tau\Z$, $T=\Im\tau$, and every $z\in\C$,
\[
 \sqrt{2T}\,|\eta_D(\tau)|^2\mathcal N(z,\tau)^2
 =\sum_{m\in\frac12\Lambda_{\mathrm{ell}}}\epsilon(m)
                         e^{-2\pi|z-m|^2/T},
\]
where $\epsilon=-1$ on $\Lambda_{\mathrm{ell}}$ and $+1$ on the three other
cosets of $\Lambda_{\mathrm{ell}}$ in $\frac12\Lambda_{\mathrm{ell}}$.
Both sides and all their derivatives converge locally uniformly.
\end{announcement}
\noindent\textit{Uses:} the theta series and Poisson summation.\par
The mechanism is comparison of Fourier coefficients on the real torus. The signed sum of the four half-period translates of the Gaussian has the same coefficients as the squared theta series with its Gaussian normalization.
\noindent\textit{Proof in the repository:} \href{https://github.com/KokunoYumeto/s6-complex-structure-workbench}{S6 complex-structure workbench}, S6\_05\_ARITHMETIC, Statement C.8.8.

We take $dd^c=(i/\pi)\partial\bar\partial$ and $dA_{z_B}=(i/2)dz_B\wedge d\bar z_B$ (ARI C.8, measure subsection). Here $z_B$ is a local base coordinate; it is unrelated to the global parameter $w$.

\begin{announcement}[ARI C.8.9; A]\label{ari:C.8.9}
On a local period chart write $\mu=a\tau+b$, with $a,b$ real,
and put $T=\Im\tau$. The positive Betti measure is
\[
 m_{P^{\mathrm{div}}}=db\wedge da
 =\frac{|\mu_{z_B}-a\tau_{z_B}|^2}{T}\,dA_{z_B}.
\]
Thus the form $da\wedge db$ in the source request labelled Lead 15(b) has the wrong sign.
The globally single-valued function $\mathcal L(t)=\log \mathcal N(P^{\mathrm{div}}_t)$ satisfies
on $\mathbb{P}^1(\C)$
\begin{equation}\label{ari:eq:ca-current}
 dd^c\mathcal L=\frac1{12}\delta_\infty-m_{P^{\mathrm{div}}},\qquad
               \int_{\mathbb{P}^1}m_{P^{\mathrm{div}}}=\frac1{12}.
\end{equation}
There are no atoms of $m_{P^{\mathrm{div}}}$ at $0$, $1$ or $\infty$.
Let $T_0=\sqrt3/2$, $\rho=(1+i\sqrt3)/2$, and
$p=\wp'_\rho((2-\rho)/3)$. The density in the ordinary
$t$ coordinate satisfies
\begin{align}
 \frac{dm_{P^{\mathrm{div}}}}{dA_t}
 &=\frac{|t|^{-2/3}}{64T_0|p|^{2/3}}
                      \bigl(1+O(|t|^{1/3})\bigr),
                         &&t\longrightarrow0,\label{ari:eq:ca-density0}\\
 \frac{dm_{P^{\mathrm{div}}}}{dA_t}
 &=\frac{|t-1|^{-3/2}}{24\pi^2\sqrt{E_4(i)}}
                      \bigl(1+O(|t-1|^{1/2})\bigr),
                         &&t\longrightarrow1.\label{ari:eq:ca-density1}
\end{align}
\end{announcement}
\noindent\textit{Uses:} the normalized theta norm; local root expansions and the cusp branch.\par
The mechanism is the curvature of the real Gaussian factor and the Jacobian of the two Betti coordinates. The cusp logarithm supplies the point mass, and the two finite root expansions supply the stated integrable densities; we retain the corrected sign and total mass from the source.
\noindent\textit{Proof in the repository:} \href{https://github.com/KokunoYumeto/s6-complex-structure-workbench}{S6 complex-structure workbench}, S6\_05\_ARITHMETIC, Statement C.8.9.

\begin{announcement}[ARI C.8.10; A; the correctly normalized equidistribution]\label{ari:C.8.10}
Let $D_n$ be the divisor, with intersection multiplicity, of
parameters $t\in\mathbb{P}^1\setminus\{0,1,\infty\}$ for which
$nP^{\mathrm{div}}_t=O$. Then
\[
 \deg D_n=\frac{n^2}{12}+O(1),\qquad
 \frac1{n^2}[D_n]\ \longrightarrow\ m_{P^{\mathrm{div}}},\qquad
 \frac1{\deg D_n}[D_n]\ \longrightarrow\ 12m_{P^{\mathrm{div}}}
\]
weakly as $n\to\infty$ through indices for which $D_n\ne0$.
For any nonrepeating sequence of individual algebraic torsion
parameters, their Galois-orbit probability measures likewise
converge to $12m_{P^{\mathrm{div}}}$. Thus the probability-limit assertion
in the source request labelled Lead 15(b) requires the factor $12$.
\end{announcement}
\noindent\textit{Uses:} ARI C.8.9; the height metric and its equidistribution theorem.\par
The mechanism applies the height metric and the small-height equidistribution theorem of DeMarco and Mavraki \cite[Theorem~1.1 and Corollary~1.2]{DM20}, with the normalization of ARI C.8.9. The passage to full division divisors uses bounded multiplicities at a fixed torsion parameter.
\noindent\textit{Proof in the repository:} \href{https://github.com/KokunoYumeto/s6-complex-structure-workbench}{S6 complex-structure workbench}, S6\_05\_ARITHMETIC, Statement C.8.10.

\begin{announcement}[ARI C.8.11; A; the two first-derivative identities]\label{ari:C.8.11}
On an ordinary modulus chart put
\[
 a=\frac{\Im\mu}{T},\qquad
 \mathcal W=\frac{\vartheta_1'(\mu,\tau)}{\vartheta_1(\mu,\tau)}
                         +2\pi ia,\qquad
 \mathcal R_\mu=\frac{d\mu}{d\tau}-a.
\]
Then
\[
 \mathcal L_\tau=\frac{\mathcal W}{8\pi i}(\mathcal W+4\pi i\mathcal R_\mu),\qquad
 \det\frac{d}{d\tau}
 \begin{pmatrix}\tau&\mu\\\mu&\beta^{\rm part}/6\end{pmatrix}
       =\frac{(\mathcal W+2\pi i\mathcal R_\mu)^2}{4\pi^2}.
\]
Moreover
\[
 \mathcal R_\mu=\frac{g_2}{6\pi i\wp'(\mu)}-\frac{\mathcal W}{2\pi i},
 \qquad
 \mathcal L_\tau=\frac{i\mathcal W(3\mathcal W\wp'(\mu)-2g_2)}
                         {24\pi\wp'(\mu)}.
\]
The critical-point factors in the source checkpoint K4 follow with their stated
normalization after multiplying by $6/\pi$.
\end{announcement}
\noindent\textit{Uses:} the theta heat equation and the Weierstrass differential equation.\par
The mechanism differentiates the logarithm of the theta quotient and its real Gaussian factor, and then differentiates the zero equation for the chosen elliptic logarithm. The root chart handles a half-period, where the displayed quotients are interpreted by continuation.
\noindent\textit{Proof in the repository:} \href{https://github.com/KokunoYumeto/s6-complex-structure-workbench}{S6 complex-structure workbench}, S6\_05\_ARITHMETIC, Statement C.8.11.

\section{The exact moduli disc}\label{sec:maximum}

\subsection{The two interval certificates}
The known abstract orbit-disc description is credited to \cite[Corollary~\texttt{p2:cor:bare-period}]{DGKXv1}; it leaves the maximum unevaluated. The following maximum and its arithmetic certification are FAM C.2.3 and PUSH E.4.

\begin{computed}[FAM C.2.3; T]\label{fam:C.2.3}
For the cusp normalization $\beta^{\rm part}+\tau\to0$,
\[
 M=-\frac3\pi\log6=-1.71100425815613411418\ldots,
 \qquad \operatorname*{arg\,max}_{B\setminus\{p_0\}}D_{\rm part}=\{1\}.
\]
The exact parameter discs are
\[
 0<|w|<216,
 \qquad 0<|u|<46656
\]
for the reference linearization $u=w^2$; an arbitrary reference constant
$\lambda_0$ multiplies the latter radius by $|\lambda_0|$.
The proof combines an analytic cusp bound, the quantitative local proof
C.2.2, and a completed outward-rounded interval cover of the remaining
compact domain.
\end{computed}
\noindent\textit{Uses:} FAM C.1.1 and C.2.2; the interval cover described below.\par

\begin{computed}[PUSH E.4.1; T]\label{push:E.4.1}
For all 19701 rational rectangles in the supplied certificate, and for
each $\tau$ in the rectangle, both zeros of $\wp_\tau$ satisfy
$\mathcal N^2<2$. The new largest certified upper bound is less than
\[
 1.999877571767410525.
\]
The 19701 rectangles, 36 local rectangles and 174 excluded rectangles
have disjoint interiors and cover exactly
$[0,1/2]\times[17/20,2]$. Every local rectangle is contained in
$|\tau-i|<3/20000$, and every excluded rectangle satisfies
$|\tau|<1$ throughout.
\end{computed}
\noindent\textit{Uses:} the supplied rational rectangle endpoints; the theta norm; Arb interval arithmetic.\par

\paragraph{Certificate block.}
The source certificate is the JSON file \path{f_interval_parallel_certificate.json} named in the FAM report. Its verifier is \path{f_interval_verify.py}; the interval calculation is \path{f_interval_cert.py}, scheduled by \path{f_interval_parallel.py}. The certificate has 19,911 leaves: 19,701 norm boxes, 36 boxes in the analytic local square, and 174 boxes strictly inside the unit circle. The root rectangle is $[0,1/2]+i[17/20,2]$, its exact area is $23/40$, and its maximum subdivision depth is 38. There are zero pending boxes. At 30 decimal digits the largest squared-norm upper endpoint is
\[
\frac{9999999982782673130994674616307871}
 {5000000000000000000000000000000000}<2.
\]
The norm boxes use a uniform complex contraction enclosure for a zero of $\wp_\tau$, followed by an outward upper bound for $\mathcal N^2$. Sign and lattice invariance give the same norm at both zeros. The verifier reconstructs the partition and checks the strict bounds and rational exclusions. The local square is covered by FAM C.2.2; the Fourier-product estimate for $T\ge2$ gives $\mathcal N<1.082$. These are the three parts of the certified proof in FAM C.2.3.

PUSH E.4.1 reuses only the rational rectangle endpoints and recomputes roots, preconditioners, radii, contraction estimates, values and tails using Arb through python-flint 0.9.0. The two programs are \path{code/e4_arb_certificate.py} (192 bits) and \path{code/e4_cover_and_local.py} (256 bits). The cover is checked at all 438 distinct vertical edges; each open vertical slab is partitioned once. The independent largest norm bound is
\[
1.999877571767410524070536478244687423578792301873142913
\ \pm9.26\cdot10^{-56},
\]
which is strictly below two. The value belongs to the enclosure of an upper bound, not to an evaluation of the maximum. The midpoint-radius representation is described in \cite[§§2--3]{Johansson2016}; the constructor and outward endpoints are documented in \cite[\texttt{arb}, \texttt{abs\_upper}, \texttt{upper}]{PythonFlint2026}.

The supplied review records a re-execution of all 19,701 FAM boxes with identical upper-bound strings (73 seconds, 12 processes), by Claude Fable 5.2 at effort high, using Python 3.13.9, sympy 1.13.1 and mpmath 1.3.0. The source review records a read of the verifier and interval routines. The PUSH certification is a producing-run computation by GPT-6 Astra, ChatGPT Pro mode; the intake did not rerun that code.

The attached material contains the TeX descriptions and the status records, but contains neither certificate file nor its executable verifier. Consequently their byte sizes, SHA-256 digests and current executions cannot be supplied from these attachments. We have not reconstructed missing certificates or claimed a new execution. This limitation concerns the supplied computational artifacts; the full duplication proof follows in the text. Its presence does not change the designated proof of record.

\subsection{Local parameters for the duplication proof}
Put $G_\square=g_2(i)=\Gamma(1/4)^8/(16\pi^2)$ and $D_*=-3\log6/\pi$ (ARI C.3.1).

\begin{proposition}[ARI C.3.1; T; fourth order at the maximum]\label{ari:C.3.1}
On either root branch at $\tau=i$, let
$h=\mu-(1-\tau)/2$ and $\delta=\tau-i$. Then
\begin{align*}
 \tau&=i-\frac{3\pi i}{2}h^2-\frac{9\pi^2i}{8}h^4+O(h^6),\\
 h^2&=\frac{2i}{3\pi}\delta+\frac1{3\pi}\delta^2+O(\delta^3),\\
 \mu&=\frac{1-\tau}{2}
  \mathbin{\pm}\left(\frac{2i\delta}{3\pi}\right)^{1/2}
          \left(1-\frac{i}{4}\delta+O(\delta^2)\right).
\end{align*}
The square root is taken holomorphically on the chosen root chart.
Furthermore
\begin{align}\label{ari:eq:cs-Li}
 \log \mathcal N={}&\tfrac12\log2-\frac\pi2|h|^2
   +\frac{3\pi^2}{8}|h|^4
   -\frac{3\pi^2}{4}|h|^2\operatorname{Re}(h^2)
   -\frac{5G_\square}{192}\operatorname{Re}(h^4)+O(|h|^6),\\
 D_{\rm part}={}&D_* -\frac2\pi|\delta|
  +\frac1\pi\bigl(|\delta|^2+|\delta|\operatorname{Im}\delta\bigr)
  +\frac{5G_\square}{72\pi^3}\operatorname{Re}(\delta^2)
  +O(|\delta|^3).\nonumber
\end{align}
In particular the leading terms of the local expansion in the source prompt are correct.
\end{proposition}
\noindent\textit{Uses:} the theta heat equation; the half-period Weierstrass equation; the cusp-normalized branch.\par
\begin{proof}
\noindent\textbf{Step 1.} Write $e(\tau)=\wp_\tau((1-\tau)/2)$ and
$\mathcal E_2=\pi^2E_2/3$. The half-period Taylor series is
\[
 \wp_\tau((1-\tau)/2+h)
 =e+(3e^2-g_2/4)h^2+(3e^3-g_2e/4)h^4+O(h^6).
\]
The exact coefficients are
$e(i)=0$, $e'(i)=iG_\square/(6\pi)$. This proves the claimed identity
$3e'(i)=(2i/\pi)G_\square/4$. For clarity its closed form follows from
$e=\pi^2(\vartheta_2^4-\vartheta_4^4)/3$ and the heat equation:
if $A=\vartheta_2(i)^4=\vartheta_4(i)^4$, then
$G_\square=4\pi^4A^2$ and $e'(i)=2\pi^3iA^2/3$.
The classical beta integral
$K(1/2)=\Gamma(1/4)^2/(4\sqrt\pi)$ and
$\vartheta_3(i)^2=2K(1/2)/\pi$ give
$G_\square=16K(1/2)^4=\Gamma(1/4)^8/(16\pi^2)$.
The beta evaluation follows by changing variable in
$\int_0^\infty(1+x^4)^{-1/2}\,dx
=\tfrac14 B(1/4,1/4)$; it is the lemniscatic complete period.


\noindent\textbf{Step 2.} Ramanujan's identities give $g_2'(i)=2iG_\square$. Differentiating
$e(-1/\tau)=\tau^2 e(\tau)$ twice at $i$ gives
$e''(i)=3i e'(i)=-G_\square/(2\pi)$. Inserting
$\delta=\alpha_2h^2+\alpha_4h^4+O(h^6)$ into the displayed half-period series
therefore gives
\[
 \alpha_2=-3\pi i/2,\qquad
 \alpha_4=\frac{g_2'\alpha_2/4-e''\alpha_2^2/2}{e'}=-9\pi^2i/8.
\]
The implicit-function theorem applies because $e'(i)\ne0$;
evenness in $h$ gives the stated remainder. Series inversion gives
the formulas for $h^2$ and $\mu$.


\noindent\textbf{Step 3.} Let $H(h,\tau)=\log(\vartheta_3(h,\tau)/\eta_D(\tau))$. The exact
identity from the half-period translation is
$\log \mathcal N=\operatorname{Re}H-\pi(\operatorname{Im}h)^2/
\operatorname{Im}\tau$. At $(0,i)$ its relevant Taylor data are
\[
 \begin{gathered}
 H_0=\tfrac12\log2,\quad H_\tau=0,\quad
 [h^2]H=-\pi/2,\quad [h^4]H=G_\square/48,\\
 \partial_\tau[h^2]H=-\pi i/4-iG_\square/(16\pi),\qquad
 H_{\tau\tau}=-G_\square/(24\pi^2).
 \end{gathered}
\]
Indeed $\partial_h^2\log\vartheta_3=-\wp(h+(1-\tau)/2)-\mathcal E_2$;
the heat equation at $h=0$ gives $H_\tau=ie/(4\pi)$, and
$\mathcal E_2'(i)=\pi i/2-iG_\square/(24\pi)$. These identities prove every entry.
After substitution $\tau=\tau(h)$, the pure $h^4$ coefficient of
$H$ is $-3\pi^2/8-5G_\square/192$. Expanding
$\operatorname{Im}\tau=1-(3\pi/2)\operatorname{Re}(h^2)+O(|h|^4)$
in the Gaussian gives \eqref{ari:eq:cs-Li}. Finally
\[
 |h|^2=\frac{2}{3\pi}|\delta|
          \left(1+\tfrac12\operatorname{Im}\delta+O(|\delta|^2)\right)
\]
and substitution into the quartic expression, followed by
multiplication by $6/\pi$, gives the formula for $D_{\rm part}$.
The source run records the script \path{Cspecial_expansion_checks.py} independently checks
these coefficients against theta evaluations; its output is retained.
\end{proof}

\subsection{The three-step duplication proof}
Throughout this proof $q_\tau=e^{2\pi i\tau}$, $T=\operatorname{Im}\tau$, and
\[
\Delta_{\rm mod}=q_\tau\prod_{n\ge1}(1-q_\tau^n)^{24},\qquad
 t=\frac{E_4^3}{1728\Delta_{\rm mod}},\qquad
 t-1=\frac{E_6^2}{1728\Delta_{\rm mod}}.
\]
This modular discriminant differs from the discriminant of the rational division model.

\begin{proposition}[ARI C.8.1; T; three uniform bounds on the elliptic norm]\label{ari:C.8.1}
For every $z\in\C$ and every $\tau$ in the standard modular
fundamental region, we have
\[
 \begin{array}{c|ccc}
 \text{restriction}&T\le2&T\le7/5&T\le11/10\\ \hline
 \mathcal N(z,\tau)&<17/10&<3/2&<143/100.
 \end{array}
\]
If $\wp_\tau(z)=0$ and $T\ge2$, then $\mathcal N(z,\tau)<541/500$.
\end{proposition}
\noindent\textit{Uses:} the theta product, Poisson summation and Fourier majorants.\par
\begin{proof}
\noindent\textbf{Step 1.} Taking absolute values in the theta series and completing the square
gives
\[
 e^{-\pi(\Im z)^2/T}|\vartheta_1(z,\tau)|
 \le\sum_{n\in\Z}
 e^{-\pi T(n+1/2+\Im z/T)^2}
 \le\sum_{n\in\Z}e^{-\pi Tn^2}.
\]
The last inequality follows from Poisson summation: the Fourier
coefficients of the periodized Gaussian are positive, so its maximum
is at an integer. The eta product and $n^2\ge1+3(n-1)$ for $n\ge1$
therefore imply, on $a\le T\le b$,
\begin{equation}\label{ari:eq:ca-sup}
 \mathcal N(z,\tau)\le
 \exp\!\left(\frac{r_a}{(1-r_a)^2}\right)
 \max\{F(a),F(b)\},\qquad
 F(T)=e^{\pi T/12}
 \left(1+\frac{2e^{-\pi T}}{1-e^{-3\pi T}}\right),
\end{equation}
where $r_a=e^{-2\pi a}$. Indeed,
$-\sum\log(1-r^n)\le r/(1-r)^2$. The function $F$ is convex,
as follows by expanding its denominator as a geometric series;
each summand is an exponential of a real linear function of $T$.
We use $6/7<\sqrt3/2$ and the following rational bounds:
\[
\begin{array}{c|@{\,}c@{\,}c@{\,}c@{\,}c@{\,}c}
T&6/7&1&11/10&7/5&2\\ \hline
F(T)&<14211/10000&<14116/10000&<14180/10000&<14782/10000&<16944/10000.
\end{array}
\]
The eta factors in \eqref{ari:eq:ca-sup} are less than $10047/10000$
for $a=6/7$ and less than $10019/10000$ for $a=1$.
Multiplying these fractions proves the asserted bounds, using
$[6/7,1]$ and $[1,2]$ separately for the first one.


\noindent\textbf{Step 2.} We specify elementary arithmetic for all the scalar estimates in
this proof. One may use $314159/100000<\pi<355/113$ and, for
$0\le x\le20$, the Taylor bounds
\begin{equation}\label{ari:eq:ca-exp}
 S_{80}(x)\le e^x\le
 S_{80}(x)+\frac{x^{81}}{81!}\frac1{1-x/82},
 \qquad S_{80}(x)=\sum_{j=0}^{80}\frac{x^j}{j!}.
\end{equation}
Taking reciprocals bounds negative exponentials. The remainder
bound follows by comparison of consecutive terms, not by numerical
sampling. These inequalities verify each displayed rational bound;
much shorter Taylor sums also suffice. Thus the decimal diagnostics
in the accompanying programs are dispensable.


\noindent\textbf{Step 3.} For the cusp we check directly the analytic step of supplied FAM C.2.3.
Choose $|\Im z|\le T/2$ by translation, put $b_0=43/100$ and
$r_0=e^{-4\pi}$. The Fourier series
\[
 \frac{\wp_\tau(z)}{\pi^2}
 =\csc^2(\pi z)-\frac13
   +8\sum_{n\ge1}\frac{nq_\tau^n}{1-q_\tau^n}
          \bigl(1-\cos(2\pi nz)\bigr)
\]
shows, if $|\Im z|\ge b_0$, that
\begin{align*}
 \left|\frac{\wp_\tau(z)}{\pi^2}+\frac13\right|
 &\le \frac1{\sinh^2(\pi b_0)}
 +8\left\{\frac{r_0}{(1-r_0)^3}
 +\frac1{2(1-r_0)}
 \left(\frac{e^{-2\pi}}{(1-e^{-2\pi})^2}
       +\frac{e^{-6\pi}}{(1-e^{-6\pi})^2}\right)\right\}
 <\frac13.
\end{align*}
The two terms in the last parentheses bound respectively
$r^n e^{2\pi n|\Im z|}$ and $r^n e^{-2\pi n|\Im z|}$.
This inequality excludes a zero. Hence a zero has
$|\Im z|<b_0$. The Jacobi product, $\log(1+x)\le x$, and
the monotonicity of $\cosh(\pi y)e^{-\pi y^2/T}$ on
$0\le y\le b_0$ give
\[
 \mathcal N(z,\tau)\le
 2\cosh(\pi b_0)e^{-\pi/3-\pi b_0^2/2}
 \exp\!\left(
 \frac{2\cosh(2\pi b_0)r_0}{1-r_0}
 +\frac{r_0^2}{1-r_0^2}\right)<\frac{541}{500}.
\]
For the stated monotonicity, the concave function
$\tanh(\pi y)-y$ is zero at zero and positive at $b_0$.
For $T\ge2$ the right side before substituting $T=2$ decreases:
the derivative of $-\pi T/6-\pi b_0^2/T$ is negative and
both product errors decrease. The two strict scalar inequalities
follow from \eqref{ari:eq:ca-exp}. This proves the cusp assertion
without using any certificate from FAM C.2.3.
\end{proof}

\begin{proposition}[ARI C.8.2; T; an explicit local maximum on a larger disc]\label{ari:C.8.2}
If $|\tau-i|\le1/50$ and $\wp_\tau(z)=0$, then
$\mathcal N(z,\tau)\le\sqrt2$, with equality exactly when $\tau=i$.
\end{proposition}
\noindent\textit{Uses:} ARI C.8.1 and C.3.1; the theta heat equation and Cauchy estimates.\par
\begin{proof}
\noindent\textbf{Step 1.} Write $\delta=\tau-i$, $h=z-(1-\tau)/2$ and
\[
 G(h,\tau)=\wp_\tau((1-\tau)/2+h),\qquad
 H(h,\tau)=\log\frac{\vartheta_3(h,\tau)}{\eta_D(\tau)}.
\]
We choose the holomorphic logarithm whose value at $(0,i)$ is
$\log\sqrt2$. The half-period theta identity gives
\begin{equation}\label{ari:eq:ca-H}
 \log \mathcal N(z,\tau)=\Re H(h,\tau)-\frac{\pi(\Im h)^2}{T}.
\end{equation}
Both $G$ and $H$ are even in $h$. On the polydisc
$|h|\le1/4$, $|\delta|\le1/10$, they satisfy
\begin{equation}\label{ari:eq:ca-GHbounds}
 |G|<30,\qquad |H|<2/3.
\end{equation}
Here are sufficient series estimates, including their margins.
Put
\[
 \mathcal S_k=\sum_{n\ge1}n^k e^{-9\pi n^2/10}\cosh(\pi n/2),
 \quad a_0=2\mathcal S_0,\quad a_1=4\pi \mathcal S_1,\quad a_2=8\pi^2\mathcal S_2.
\]
Then $a_0<3/10$, $a_1<19/10$, $a_2<12$.
The first two summands suffice to verify these inequalities: for
$n\ge3$, each of the two exponentials in the hyperbolic cosine
has successive-term ratio at most $e^{-29\pi/5}$, and the
factor $n^k$ increases this ratio by at most $(4/3)^2$.
We use \eqref{ari:eq:ca-exp} for the first terms and this geometric
bound for the tail. With $r=e^{-9\pi/5}$, the identities
\[
 G=-\partial_h^2\log\vartheta_3-\frac{\pi^2}{3}E_2,
 \qquad
 |E_2|\le1+\frac{24r}{(1-r)^3}
\]
give
\[
 |G|\le\frac{a_2}{1-a_0}+\frac{a_1^2}{(1-a_0)^2}
       +\frac{\pi^2}{3}\left(1+\frac{24r}{(1-r)^3}\right)<30.
\]
Also
\[
 |H|\le-\log(1-a_0)+\frac{11\pi}{120}
                    +\frac{r}{(1-r)^2}<\frac23.
\]
The logarithm is well defined since $|\vartheta_3-1|<3/10$;
the displayed eta product fixes its branch. A geometric bound
for $-\log(1-x)=\sum_{n\ge1}x^n/n$ verifies the last estimate.


\noindent\textbf{Step 2.} Let
\[
 e(\tau)=\frac{\pi^2}{3}(\vartheta_2(0,\tau)^4-
                              \vartheta_4(0,\tau)^4),\qquad
 \mathcal E_2(\tau)=\frac{\pi^2}{3}E_2(\tau),\qquad
 g=\frac{g_2(i)}4.
\]
On $|\delta|\le1/50$ the following loose bounds hold:
\begin{equation}\label{ari:eq:ca-derivatives}
 |e'|<14,\quad |e''|<60,\quad |g_2'|<500,
 \quad |\mathcal E_2'|<2,\quad |E_6''|<100,
 \qquad 47<g<95/2.
\end{equation}
For completeness these are scalar Fourier-series bounds, not
interval computations. At $T\ge49/50$, upper bounds for
$|\vartheta_j|,|\vartheta_j'|,|\vartheta_j''|$ are
\[
\begin{array}{c|ccc}
j&0\text{ derivatives}&1\text{ derivative}&2\text{ derivatives}\\ \hline
2&929/1000&742/1000&670/1000\\
4&1093/1000&290/1000&911/1000.
\end{array}
\]
They follow by differentiating the absolutely convergent theta
series, taking absolute values, retaining two terms, and using
the same successive-term bound. Differentiating the fourth
powers proves the first two bounds in \eqref{ari:eq:ca-derivatives}.
For $r\le1/470$, the remaining derivative majorants are
\begin{align*}
 |g_2'|&\le640\pi^5\frac{r(1+11r+11r^2+r^3)}{(1-r)^7}<500,\\
 |\mathcal E_2'|&\le16\pi^3\frac{r(1+r)}{(1-r)^5}<2,\\
 |E_6''|&\le2016\pi^2
 \frac{r(1+r)(1+120r+1191r^2+2416r^3+1191r^4+120r^5+r^6)}
 {(1-r)^{11}}<100.
\end{align*}
These follow by summing $\sum n^k r^n$ and bounding
$(1-r^n)^{-1}$ by $(1-r)^{-1}$. Substitution of $r=1/470$
and $\pi<355/113$ verifies the rational inequalities.
Finally $g=\pi^4E_4(i)/3$ and
$E_4=1+240\sum n^3q_\tau^n/(1-q_\tau^n)$ give $47<g<95/2$.
The identity of ARI C.3.1, checked by differentiating $e$ and using
the theta heat equation, reads
\[
 e(i)=0,\qquad e'(i)=id,\qquad d=\frac{2g}{3\pi}>99/10.
\]


\noindent\textbf{Step 3.} The differential equation $\wp''=6\wp^2-g_2/2$ determines the
first coefficients in the even expansion:
\[
 G(h,\tau)=e+a h^2+ea h^4+\mathcal R_G,
 \qquad a=3e^2-g_2/4.
\]
Consequently, on $|\delta|\le1/50$,
\[
 |e|\le7/25,\quad |a+g|\le\frac52+3(7/25)^2=\frac{3419}{1250},
 \quad |a|\le\frac{62794}{1250}.
\]
Cauchy's estimate from \eqref{ari:eq:ca-GHbounds} yields
\[
 |\mathcal R_G|\le
 \frac{30(4|h|)^6}{1-16|h|^2}.
\]
On $|h|=1/10$ the term $ah^2$ has modulus at least
$55331/125000>0.4426$, whereas the sum of the other terms is at most
\[
 \frac7{25}+\frac7{25}\frac{62794}{1250}\frac1{10000}
       +\frac{30(4/25)^3}{1-4/25}<0.4277.
\]
Rouch\'e's theorem puts two zeros, with multiplicity, in
$|h|<1/10$. The disc injects into the torus and $\wp$ has exactly
two zeros on the torus, so it contains both of them.


\noindent\textbf{Step 4.} Write $G=e+h^2Q$. The preceding estimates give, for $|h|\le1/10$,
\[
 |Q+g|\le
 \frac{3419}{1250}+\frac7{25}\frac{62794}{1250}\frac1{100}
       +\frac{30(4/25)^3}{(1-4/25)/100}<\frac{88}{5}.
\]
Taylor's formula and $|e''|<60$ give
$|e|\ge(99/10-3/5)|\delta|=(93/10)|\delta|$.
At either zero, therefore,
\begin{equation}\label{ari:eq:ca-delta-h}
 |\delta|<7|h|^2,
\end{equation}
because $(95/2+88/5)/(93/10)=7$.


\noindent\textbf{Step 5.} Expand
$H=H_0(\tau)+H_2(\tau)h^2+H_4(\tau)h^4+\mathcal R_H$.
The heat equation and the identity for $G$ give
\[
 H_0'=\frac{ie}{4\pi},\quad
 H_2=-\frac{e+\mathcal E_2}{2},\quad H_4=-\frac a{12}.
\]
Thus $H_0(i)=\log\sqrt2$, $H_0'(i)=0$,
$|H_0''|<6/5$, $H_2(i)=-\pi/2$, $|H_2'|<8$, and
$|H_4|<21/5$. Cauchy's estimate gives
\[
 |\mathcal R_H|\le
 \frac{(2/3)(4|h|)^6}{1-16|h|^2}
 \le\frac{2048}{63}|h|^4
 \quad (|h|\le1/10).
\]
Combining \eqref{ari:eq:ca-H} and \eqref{ari:eq:ca-delta-h}, we obtain
\begin{align*}
 \log \mathcal N-\log\sqrt2
 &\le-\frac\pi2\Re(h^2)-\pi(\Im h)^2+C|h|^4\\
 &=-\frac\pi2|h|^2+C|h|^4,\\
 C&=\frac{147}{5}+56+\frac{21}{5}
                 +\frac{2048}{63}+\frac{1100}{49}<150.
\end{align*}
The final term of $C$ bounds the change of $1/T$ from $1$:
$7\pi/(1-1/50)<1100/49$. Since $150|h|^2<3/2<\pi/2$,
the result is strict unless $h=0$. Equation
\eqref{ari:eq:ca-delta-h} then gives $\delta=0$, where equality
holds by the checked isogeny computation at $i$.
\end{proof}

\begin{proposition}[ARI C.8.3; T; the comparison of the two local discs]\label{ari:C.8.3}
Every complex number $t_0$ with $|t_0-1|\le1/500$ has a
representative $\tau$ with $|\tau-i|<1/50$ and $j(\tau)=1728t_0$.
Consequently every zero of $\wp$ for that modulus satisfies
$\mathcal N\le\sqrt2$, with equality only for $t_0=1$.
\end{proposition}
\noindent\textit{Uses:} ARI C.8.2; the Ramanujan differential identities and Rouche theorem.\par
\begin{proof}
\noindent\textbf{Step 1.} At $i$, $E_6=0$ and the Ramanujan differential identity gives
$|E_6'(i)|=\pi E_4(i)^2>13/2$, since $E_4(i)>29/20$.
On $|\delta|=1/50$, \eqref{ari:eq:ca-derivatives} gives
\[
 |E_6(\tau)|>\frac{13}{100}-\frac1{50}=\frac{11}{100}.
\]
On the whole disc, $|q_\tau|<1/470$ and
$|\Delta_{\mathrm{mod}}|\le |q_\tau|\exp(24|q_\tau|/(1-|q_\tau|))<1/400$.
It follows on the boundary that
\[
 |t-1|=\frac{|E_6|^2}{1728|\Delta_{\mathrm{mod}}|}
 >\frac{121}{43200}>\frac1{500}.
\]
Rouch\'e applied to $E_6'(i)\delta$ shows that $E_6$ has
exactly one zero in the disc. Hence $t-1$ has two zeros there,
with multiplicity, and another application of Rouch\'e shows
that $t-t_0$ has two for each $|t_0-1|\le1/500$.
This argument proves coverage of the $t$-disc, rather than
inferring coverage merely from a boundary bound. Elliptic
curves over $\C$ with the same $j$ are isomorphic. In terms of
lattices the isomorphism is a homothety; it carries the two
zeros of $\wp$ to the two zeros of the scaled $\wp$ and
preserves $\mathcal N$. Statement C.8.2 applies to either zero.
\end{proof}

\begin{proposition}[ARI C.8.4; T; the first three normalized duplication steps]\label{ari:C.8.4}
Put $v_t=t-1$. For $n=2^k$, $k=1,2,3$, and $nP^{\mathrm{div}}\ne O$,
\begin{equation}\label{ari:eq:ca-ladder}
 \left(\frac{\mathcal N(P^{\mathrm{div}})}{\sqrt2}\right)^{n^2}
 =\frac{3^{3/4}}{2n}\,
       \frac{\mathcal N(nP^{\mathrm{div}})|v_t|^{-1/4}}{|\mathcal A_k(v_t)|},
\end{equation}
where
\begin{align*}
 \mathcal A_1(v_t)&=1,\qquad \mathcal A_2(v_t)=1-5v_t/27,\\
 \mathcal A_3(v_t)&=1-\frac{217}{27}v_t-\frac{20486}{729}v_t^2
       +\frac{10910}{19683}v_t^3
       +\frac{161383}{177147}v_t^4
       -\frac{245}{14348907}v_t^5.
\end{align*}
On $|v_t|\le1/28$ we have $|\mathcal A_3(v_t)|\ge1-10|v_t|$.
For every $k\ge1$, not just these first three steps, define
\[
 E_k(t)=\frac{|\Delta_E|^{(4^k-1)/12}|t-1|^{1/4}}
 {2^{4^k/2}|\psi_{2^k}(P^{\mathrm{div}})|}.
\]
Then $(\mathcal N(P^{\mathrm{div}})/\sqrt2)^{4^k}=E_k(t)\mathcal N(2^kP^{\mathrm{div}})|t-1|^{-1/4}$,
and the removable limiting value is
\[
 \lim_{t\to1}E_k(t)=\frac{3^{3/4}}{2^{k+1}}
                 =2^{1-k}\frac{3^{3/4}}4.
\]
\end{proposition}
\noindent\textit{Uses:} ARI C.4.2, C.4.4 and C.3.1.\par
\begin{proof}
\noindent\textbf{Step 1.} Substitute the exact division polynomials from Statement C.4.2
in \eqref{ari:eq:cs-division}. Equivalently, remove
$t^{2(n^2-1)/3}(t-1)^{n^2/4}$ from $\psi_n(P^{\mathrm{div}})$ and divide
the resulting polynomial by its value at $t=1$. For $n=2,4,8$
this gives the displayed polynomials and the constant
$3^{3/4}/(2n)$. All operations are rational polynomial operations;
\path{code/C_ladder_algebra.py} reproduces them and records
the next two steps as well. The general identity is C.4.2.
For its limit write $\mu=(1-\tau)/2+h$ as in C.3.1. Since
$n=2^k$ is even, $n\mu$ is congruent to $nh$ modulo the
period lattice. Hence
\[
 \mathcal N(nP^{\mathrm{div}})=2\pi|\eta_D(i)|^2 n|h|(1+o(1)),\qquad
 t-1=\tfrac94\pi^4E_4(i)h^4(1+o(1)).
\]
The identity $E_4(i)^3=1728\eta_D(i)^{24}$ then gives
$\mathcal N(nP^{\mathrm{div}})|t-1|^{-1/4}\to2n/3^{3/4}$.
Since the left side of the general ladder tends to one,
this proves the stated limit for every $k$.
The triangle inequality gives
$|\mathcal A_3(v_t)|\ge1-|v_t|M_{\rm coeff}$, where
\[
 M_{\rm coeff}\le\frac{81}{10}+\frac{29}{28}
                +\frac1{28^2}+\frac1{28^3}+\frac1{28^4}<10.
\]
Here we bounded respectively the five absolute coefficients by
$81/10,29,1,1,1$. This is a rational inequality on one disc,
not a subdivision or a root search.
\end{proof}

\begin{proposition}[ARI C.8.5; T; FAM C.2.3 by explicit analytic estimates]\label{ari:C.8.5}
For every lattice $\Z+\tau\Z$ and every zero $z$ of its
Weierstrass function,
\[
 \mathcal N(z,\tau)\le\sqrt2.
\]
Equality holds exactly for the modular orbit of $i$.
In the parameter of the family this says
$\max D_{\mathrm{part}}=-(3/\pi)\log6$, only at $t=1$.
\end{proposition}
\noindent\textit{Uses:} ARI C.8.1--C.8.4; the norm values of ARI C.1.2.\par
\begin{proof}
\noindent\textbf{Step 1.} Reduce to the standard modular region. The portion $T\ge2$
is covered by Statement C.8.1. Suppose $T\le2$ and put
$r=|t-1|$. For $r\ge9/10$, the step $n=2$ in
\eqref{ari:eq:ca-ladder} and $\mathcal N(2P^{\mathrm{div}})<17/10$ give strict inequality,
because
\[
 \frac{27(17/10)^4}{256}<\frac9{10}.
\]
There is no $2$-torsion point of the section on the smooth base.


\noindent\textbf{Step 2.} If $r\le9/10$, then $T<7/5$. Indeed, $T\ge7/5$ would give
$|q_\tau|<1/6500$ and
\[
 |E_4|>24/25,\qquad |\Delta_{\mathrm{mod}}|<1/6400,
 \qquad |t|>\frac{6400(24/25)^3}{1728}>19/10,
\]
contrary to $|t|\le19/10$. The estimates follow from
\[
 \sum_{n\ge1}\sigma_3(n)r^n
 \le\frac{r(1+4r+r^2)}{(1-r)^5},\qquad
 |\Delta_{\mathrm{mod}}|\le r\exp\frac{24r}{1-r}.
\]
For $1/28\le r\le9/10$, the $n=4$ step and
$\mathcal N(4P^{\mathrm{div}})<3/2$ now apply. The function
$(1-5r/27)r^{1/4}$ is increasing up to $r=27/25$.
At its left endpoint its fourth power is
\[
 \frac1{28}\left(\frac{751}{756}\right)^4
 >\frac{2187}{65536}
 =\left(\frac{3^{3/4}}8\frac32\right)^4.
\]
This proves strict inequality on the whole annulus. The lower
bound for $\mathcal A_2$ is positive there, so no division by a zero has
occurred.


\noindent\textbf{Step 3.} If $r\le1/28$, then $T<11/10$. For otherwise $|q_\tau|<1/1000$,
and the Lambert series with
\[
 \sum_{n\ge1}\sigma_5(n)r^n
 \le\frac{r(1+26r+66r^2+26r^3+r^4)}{(1-r)^7}
\]
gives $|E_6|>9/20$, $|\Delta_{\mathrm{mod}}|<1/900$, and hence
$|t-1|>900(9/20)^2/1728>1/28$.
On $1/500\le r\le1/28$, the $n=8$ step,
$\mathcal N(8P^{\mathrm{div}})<143/100$, and $|\mathcal A_3|\ge1-10r$ prove the claim.
Indeed $(1-10r)r^{1/4}$ has a single critical point in this
interval, a maximum at $r=1/50$, so its minimum is at an endpoint.
At both endpoints its fourth power exceeds $9/5000$, whereas
\[
 \left(\frac{3^{3/4}}{16}\frac{143}{100}\right)^4
 <\frac7{4000}<\frac9{5000}.
\]
Again the displayed positive bound excludes a division-polynomial
zero. The remaining disc $r\le1/500$ is covered by Statement
C.8.3. The smooth limiting fibre at $t=0$ has
$\mathcal N=3^{1/4}<\sqrt2$; the formula also extends there by
continuity, so the singular rational model causes no omitted
modulus. The only equality occurs at $t=1$.
\end{proof}

\subsection{Announcement on rays and level curves}

\begin{announcement}[ARI C.8.15; A; every critical point and every level set]\label{ari:C.8.15}
For every zero $\mu$ of $\wp_\tau$, the zero sets of the two
critical factors are exactly
\[
 \begin{aligned}
 \mathcal W=0&\quad\Longleftrightarrow\quad
       \tau\in\operatorname{SL}_2(\Z)\{i,\rho\},\\
 3\mathcal W\wp'_\tau(\mu)-2g_2=0&\quad\Longleftrightarrow\quad
       \tau\in\operatorname{SL}_2(\Z)\rho.
 \end{aligned}
\]
On the coarse $t$-plane the function $\mathcal L$ has no critical point
away from $t=1$. The point $t=0$ is a noncritical $C^1$ point,
with $\mathcal L_t(0)=7/384$. The point $t=1$ is the unique global
maximum, of value $\log\sqrt2$, and has the fractional-power
local maximum described in C.3. In particular there is no
saddle in $0<t<1$. On the root charts, $h=0$ over $t=1$
is a nondegenerate maximum, while $\delta=0$ over $t=0$
is a degenerate critical point with nonzero harmonic cubic
leading term and three alternating ascending and descending
sectors. These are all critical points on the root charts.

For every angle $\theta$, $r\mapsto \mathcal L(1+re^{i\theta})$ is
strictly decreasing for $r>0$, from $\log\sqrt2$ to $-\infty$.
For each $c<\log\sqrt2$, the level set $\{\mathcal L=c\}$ is one
Jordan curve, strictly star-shaped about $1$. Its bounded
complement is $\{\mathcal L>c\}$; its unbounded complement is
$\{\mathcal L<c\}$. The curve is real analytic except possibly at
$t=0$ when $c=\frac14\log3$, where it is $C^1$ with nonzero
tangent. The closed superlevel sets form nested topological
discs. At the maximum the level set is the single point $1$,
and above it the level set is empty. The corresponding
statements for $D_{\rm part}$ follow by its increasing affine
change from $\mathcal L$.
\end{announcement}
\noindent\textit{Uses:} ARI C.8.11--C.8.14 and the two local root expansions.\par
The mechanism bounds the two critical factors on a cusp region and on the remaining compact modular region. The root expansions identify their behavior over $t=0,1$; the radial derivative then gives the one-radius parametrization of each lower level.
\noindent\textit{Proof in the repository:} \href{https://github.com/KokunoYumeto/s6-complex-structure-workbench}{S6 complex-structure workbench}, S6\_05\_ARITHMETIC, Statement C.8.15.
